% =========================================================================
\chapter{Reproduction and Baseline Diagnosis}
\label{ch:baseline}
% =========================================================================



\section{Mechanisms in the released loop and differences from the paper}
\label{sec:official-strengths}
% -------------------------------------------------------------------------

Four implemented mechanisms are relevant to the following diagnosis.

\begin{enumerate}
  \item \textbf{Success-trajectory analysis.} The all-pass Digester
        template instructs the role to identify a reusable strategy and any
        latent fragility.
  \item \textbf{Selective invocation and access scope.} Planning and
        evolution are conditional (\S\ref{sec:fielded}): whether they run
        in a round depends on an actionability score computed mechanically
        from the dossier mix. Scope gates are configured to restrict access to the harness
        source and other runs' histories, and to reject absolute paths under
        those roots in shell commands. These are the restrictions
        \S\ref{sec:governance} reports tightening after an adjudication
        file was left readable in a run root.
  \item \textbf{Ordered candidate validation.} Five deterministic gates
        (\S\ref{sec:motivation}) run in sequence. The structure gate's
        twelve evidence invariants and the counterfactual replay provide the
        basis for this thesis's sixth gate.
  \item \textbf{Prediction-based feedback and the rollback clamp.} Each
        shipped candidate names the tasks it expects to flip; the next round
        scores that prediction, records its hit rate, and lists in the
        Planner's summary any predicted task that still fails. The only
        automatic rollback is the post-round score clamp of
        \S\ref{sec:open-loop}, keyed to the total score rather than to the
        prediction. This prediction record informs the localization metric
        of \S\ref{sec:localization}; \S\ref{sec:open-loop} reports whether
        the clamp engaged.
\end{enumerate}

\subsection*{Where the release diverges from the paper}

Five specifications in the paper are absent from the code: the variant
pool and ensemble routing behind the GAIA headline (the release is one
lineage with a \texttt{ship\,|\,no\_op} decision); the global ratchet
forbidding a solved task to be given back; the $\pm 5\%$ noise threshold on
retention; three seeds per cell, the code defaulting to $k{=}1$; and the
Planner's per-candidate briefs. Appendix~\ref{app:deviations} records each
with its code pin.

The measures taken are the same in every case. The tree is vendored
byte-for-byte under a hash manifest, so what was run is the release and not
a repair of it; the authors' own upstream fix is adopted where they
published one; the missing specifications are supplied by the
\emph{experiment} rather than by patching the system --- three seeds per arm, the paper's $\pm5\%$ convention read at the analysis
layer, an empirical same-configuration envelope under the declared post-data
window ruling, and offline regression accounts --- and campaign outcome statistics use archived task histories or trajectories
under the ledger's declared reading conventions. Publication transcriptions,
code inspections, logs and inventories are separately classed in the ledger. Where a finding depends on one of these
divergences, the dependence is stated.

% -------------------------------------------------------------------------
\section{System execution and study scope}
\label{sec:fielded}
% -------------------------------------------------------------------------

The paper describes a round as execution of the current harness on a task
batch, digestion of the resulting trajectories, conditional planning and
evolution, and candidate review by a Critic followed by deterministic checks
before commit~\cite{aegis}. The released study configuration advances one lineage from round $k$ to
round $k+1$. Each round is bounded by a token and cost budget and by a cap on the
number of tasks attempted. Differences between the paper and release are
listed in \S\ref{sec:official-strengths}; the rest of this chapter concerns
the released implementation.

% -------------------------------------------------------------------------
\section{Open-loop hazards: feedback written but not delivered}
\label{sec:open-loop}
% -------------------------------------------------------------------------

A self-evolving loop is only closed if the consequences of an edit reach
the author of the next edit. Several feedback paths break that delivery:
regression and verdict paths here; the candidate-rejection ledger and
Planner targeting in \S\ref{sec:decision}. The score clamp and missing
ratchet are retention controls, not delivery channels. These are properties
of the released tree, not of campaign execution or the configured model tier.

\paragraph{The regression report is left empty by a round-indexing slip.}
Each round the loop computes which previously-passing tasks have started
failing --- the signal a lineage without a ratchet most needs. An
off-by-one in the round index draws the comparison against an empty set,
so the report is structurally always empty. The authors' own upstream fix
is adopted here and every regression account is re-derived offline.

\paragraph{The verdict contract is stated twice, and the two statements
have drifted apart.} The schema a role is instructed to emit and the schema
the consumer validates against no longer agree, so validation fails
softly: the outcome is dropped rather than raised, and the next author is
never told what happened to the last edit.

\paragraph{The rollback threshold cannot engage on a small bed.} Rollback
fires only when the pass \emph{rate} drops by at least $0.05$ \emph{and}
the pass \emph{count} drops by at least $3$ against the last validated
round (\texttt{run\_meta\_aegis.py:1425}). The conjunction is a property
of the code: wherever fewer than three tasks were passing, the count
condition cannot be met at all, and where exactly three were, only a
total loss meets it. The clamp is arithmetically unreachable in exactly
the situation --- new wiring, first exercised on a handful of tasks ---
for which a clamp exists.

\paragraph{The floor the paper promises is absent: a single lineage with
no global ratchet.} This is not a delivery channel but the floor beneath
them. The paper specifies a global monotone ratchet: no previously solved task may be
given back. The code implements no such thing. What it has instead is an
advisory watchlist over adjacent rounds, which the agent may overrule by
declaring a regression transient; a counterfactual replay gate; and the
post-round score clamp of the preceding paragraph. None of these is a floor. In flight the watchlist fires at the noise rate --- hard regressions
before 81 of 86 evolve rounds across the six whitelist campaigns --- and
the clamp fired five times, reverting seven landed candidates on drops of
5 to 9 tasks, three of them inside the same-config band --- the
$-8\ldots{+}5$ union of \S\ref{sec:power}, used throughout as a descriptive
reference range \F{66}\F{15}.

The tree guards its own worst known regression --- an Evolver edit that
replaced the system-prompt builder and thereby deleted an injected domain
policy, recorded in the authors' own words as a ``catastrophic
regression'' --- with an instruction inside a prompt template rather than
a code path, and the only committed performance record in the repository
is a commit message recording that lineages \emph{never recovered} from
one such crash \F{19}. This unrecovered lineage failure motivates attaching a readout
to the retention decision (\S\ref{sec:future}).

% -------------------------------------------------------------------------
\section{The measurement layer: three contamination sources}
\label{sec:measurement}
% -------------------------------------------------------------------------

\paragraph{Judge false positives can flip the sign of a round.} Scoring is
performed by an LLM judge, named with the other meta-role models in
\S\ref{sec:beds}. The judge is \emph{structurally hardened}: it must quote
the sentence where the agent commits to an answer, and what it cannot quote
fails, so silence and reasoning that reaches no commitment cannot pass.
Without that requirement a recorded $40\%$ of answerless trajectories
passed, on an unrecorded sample; with it, five of five known cases were
graded correctly, and the error rate on these six campaigns was never
measured \F{69} (\S\ref{sec:threats}). Every between-arm comparison in this
thesis uses audited scores only \F{46}.

\paragraph{The bed's ground truth is partly on the public web.} Of a
49-task cohort assembled from the whitelist graph campaign's swing tasks,
thirteen are
\emph{leak-route}: more than half of their passes touch benchmark
artifacts, and reference answers appear verbatim inside solver
trajectories \F{31}. 

\paragraph{Re-measuring the same configuration moves 20\% of the bed.}
This finding governs the design of every experiment in this thesis; the
arithmetic is worked in \S\ref{sec:power}. Across the six ship-free windows of this
campaign's sixteen-round readout, $21.0\%$ of task outcomes flip between
two consecutive rounds running an identical configuration --- $126$ of
$600$ task-pairs --- and the resulting total-score swing runs from $-4$ to
$+5$ with no treatment applied at all. Over all twenty ship-free windows of
the three no-graph seeds that swing runs from $-8$ to $+5$, and it is the
union, not this campaign's own range, that the thesis uses as its
descriptive reference range; \S\ref{sec:power} derives it and \F{15} records the
per-seed ranges.

The third source is not a defect hardening removes: it is a property of
the instrument and replicates across three seeds
(\S\ref{sec:replication}).

% -------------------------------------------------------------------------
\section{The evidence-pipeline audit}
\label{sec:evidence-audit}
% -------------------------------------------------------------------------

Every dossier the Digester produces carries, for each tool call, a verdict
on whether the call's output was \emph{used}. That verdict is the loop's
only signal about whether an action mattered. It is consumed by the
Planner when ranking directions, quoted by the Evolver when justifying a
candidate, and reviewed by the Critic when deciding whether to veto. Across the available first-repetition traces from 14 task rounds of the
baseline campaign's 103-task trajectory set, the column judges only
$8.1\%$ of payload calls to have been used \F{1}. The defect is in the
column, and is established in three tiers of decreasing strength; the
argument needs only the first two, the third resting on a premise the data
cannot supply.

\subsection*{(i) Construction}

The column tests whether the \emph{immediately next} assistant message
shares a verbatim substring of at least twenty characters with the tool
output (\texttt{harnessx/\allowbreak aegis/\allowbreak stages/\allowbreak trace\_facts.py}). The rule cannot identify paraphrased or delayed use as such; incidental
qualifying overlap in the immediately next message could still produce a
positive verdict. When that message has no overlap of at least twenty
characters with the tool output, the verdict is necessarily ``not
referenced'', even if shorter spans were genuinely reused.

\subsection*{(ii) Exposure on the selected diagnostic subset}

Among the 419 overlap-rule-negative answer-carrying calls, the
reference-answer median is \textbf{9} characters \F{3}. This selected subset
makes the threshold mismatch salient; it does not estimate its bed-level
prevalence or the rule's false-negative rate.



\subsection*{(iii) Rate, and what it does not license}

The rate measures the reach of the defect. Of the $559$ calls whose output literally
contains the final answer of a rollout that passed, \textbf{419
($75.0\%$)} receive the overlap rule's ``not referenced'' label --- by
coincidence the same $75.0\%$ as the column's NO rate over all payload
calls \F{1}, a different population --- and
$96.7\%$ of those carry an answer below the column's own minimum length
\F{3}. The figure is robust
across variants of the answer-length window at $73.8$--$75.1\%$ \F{2};
the column's own twenty-character threshold was not swept, and by tier (i)
it could not have been swept without changing the instrument.

The limit: a call whose output contains the answer, on a path that passed,
is a call sitting on the successful path holding the answer, which is not
the same as having caused the pass. The finding
therefore rests on (i) and (ii), which need no causal premise, and the
$75.0\%$ is reported as a property of the column's output --- the rate at
which it applies its overlap verdict to answer-carrying calls --- not as a
count of verified causal errors \F{2}.

\subsection*{The disclaimer, and why it makes the case worse}

The Digester's own text does carry a caveat. It reads:

\begin{quote}\small
\emph{Heuristic:} assistant message after the tool call does not contain
any $\ge$20-char substring from the tool output. This is the single most
common way a tool call ``fires but fails to affect the model'' --- e.g.\
when protocol flattening drops multimodal content.
\end{quote}

The disclaimer names one failure mode --- multimodal flattening --- and
never mentions that it cannot credit \emph{verbatim reuse confined to
spans shorter than twenty characters}. It describes the column as detecting calls that fire without
affecting the model, which is precisely the interpretation tier (ii)
forbids. And it does not travel with the verdict: the per-call rows in the
dossier's main table carry a bare \textbf{NO}.

\subsection*{Reach into decisions}

The verdict propagates. It is recorded in \textbf{86.0\%} of available dossiers \F{4}, it
is quoted verbatim inside three candidate documents \F{5}, and one
candidate was shipped on it \F{6}.

That shipped candidate named three tasks it expected to flip and hit none
of them. \emph{This is reported as illustration, not as inference.} Against
the round's base rate of $7/32$, a candidate with no effect whatever
returns zero of three with probability $0.48$ \F{6}. The specimen shows how the
column propagated into a shipped decision; it says nothing about
whether that decision was good. The same evidentiary standard that keeps the
no-graph localization lift of \S\ref{sec:localization} a descriptive
selected-candidate association applies with more force at $n{=}1$, and is
applied here.

% -------------------------------------------------------------------------
\section{Decision-plane pathologies}
\label{sec:decision}
% -------------------------------------------------------------------------

Three further properties of the decision plane were recorded; all are
code-pinned.

\paragraph{Free-text edit identity is not verified downstream.} The Evolver
composes a manifest and a \texttt{config.yaml} as prose. The edit surface can
install a new processor class (\S\ref{sec:one-repair}), but the release has
no downstream check that the candidate document and applied change describe
the same intervention.

\paragraph{The candidate-rejection ledger is written and never read.} The loop
records why each candidate was rejected; nothing consumes it: across the six
campaigns 94 candidates were rejected and 3 revived \F{67}. This is another broken feedback path described in
\S\ref{sec:open-loop}.

\paragraph{The aimer is never shown where its aims landed.} The guidance
rebind list omits the Planner (\texttt{planner.py:30--46}), in every arm,
so the role that decides what to aim at never sees where previous aims
landed: whatever the loop learns about which targets convert does not reach
the stage that chooses targets.

% -------------------------------------------------------------------------
\section{The mechanism-consistent guard event}
\label{sec:one-repair}
% -------------------------------------------------------------------------

One post-hoc baseline event has a mechanism-readable score movement. It is
mechanism-consistent evidence, not an isolated estimate of one intervention's
effect.

After its first round the loop diagnosed a shell-quoting bug in the
harness's Bash tool, authored a Windows bash guard and shipped it into R2;
budget-exhaustion exits did not move ($36 \to 35$) \F{39}. Its next proposals
shipped into R3: a step-countdown processor that tells the solver how many
steps remain, a prompt and a PDF tool. In that round
\texttt{budget\_exceeded} exits collapsed from $35$ to $4$; twelve of the
sixteen tasks that newly passed had exited after budget exhaustion in the
preceding round; the GAIA difficulty tier 2 subset gained $16.3$ percentage
points; and the total score rose by $9$ tasks, clearing the same-config band
\F{39}. Three candidates shipped into that round; without a same-window
comparison their individual effects are not separated. It is the largest
single-round move that campaign makes after its baseline draw, and --- of
the six above-band whole-bed moves in all six campaigns --- the only one
whose mechanism is legible in the exit ledger \F{48}.

The two calibers of same-configuration noise disagree, and not in this
thesis's favour: the band is a range over twenty windows, not a standard
deviation, and the paired-change SD of $3.70$ in \S\ref{sec:power} is not
a score-level sigma. The repair stands on the
mechanism counts under either caliber; how small an effect the score axis
can register remains unresolved.

\paragraph{How to read it.} The adjacent-round observation is consistent
with repair of an engineering defect: budget-exhaustion exits fell while the
score rose. Three candidates shipped concurrently and there is no
same-window comparison, so the observation does not isolate the processor's
efficacy or establish a capability change. Nor does it record retention:
the processor stayed for four rounds and left at R7, the round a
three-bucket candidate shipped, after which exits ran 6--12 per round,
while the earlier bash guard stayed to R15. Chapter~\ref{ch:results}
returns to that record and its cross-campaign scope \F{39}, and separately
to the score response (\S\ref{sec:gain-face}).
